\documentclass[12pt, a4paper, bibliography=totoc]{scrartcl}


%\input{~/Latex/article/article_config.tex}
% personal data
\date{\today}


% language
\usepackage{polyglossia}
\setmainlanguage{english}
\setotherlanguages{german}
\usepackage{microtype}
\usepackage{dcolumn}

\usepackage[style=numeric,
			natbib=true,
			backend=biber]{biblatex}		%Bibliographie
\usepackage[autostyle=true,
			 german=quotes]
			 {csquotes}					%Anführungszeichen
\usepackage{blindtext}


%math and theorems
\usepackage{amsmath}
\usepackage{amssymb}
\usepackage{amsopn}					%Matheoperatoren
\usepackage[amsmath,thmmarks,hyperref]{ntheorem}
\usepackage{mathtools}
\usepackage{mathdots}					%Punkte
\usepackage{dsfont}
\usepackage{upgreek}					%Griechische Buchstaben
\usepackage{bbm}						%Mengensymbol
\usepackage{physics}					%Physiksymbole
\usepackage{relsize}						%Größenangaben
\usepackage[separate-uncertainty,
			per-mode=symbol]
			{siunitx}					%Einheiten
%\usepackage{tikz}						%Zeichnen
\usepackage{upgreek}					%Griechische Buchstaben
\usepackage{enumitem}
\setlist{nolistsep}


%useful packages
%\usepackage{geometry}
\usepackage{xcolor}
\usepackage{graphicx}
\usepackage{float}
\usepackage{csquotes}
\usepackage{todonotes}
\usepackage{booktabs}
\usepackage{array}
\usepackage[labelfont=bf]{caption}
\usepackage{wrapfig}
\usepackage{enumitem}
%\usepackage{xr} % cross referencing
%\usepackage{titling}
%\usepackage{titlesec}
%\usepackage[Bjornstrup]
%			{fncychap}					%Kapitellayout


\setmainfont{Linux Libertine O}
\setsansfont{Linux Biolinum O}

\usepackage{scrhack}					%Verbesserung Pakete
\usepackage{xltxtra}						%fontec


\newcommand{\im}{\mathrm{i}}
\newcommand{\e}{\mathrm{e}}
\renewcommand{\pi}{\uppi}
\renewcommand{\epsilon}{\varepsilon}


\addbibresource{bibliography.bib}

%color settings
\definecolor{myred}{RGB}{196,19,47} 
\definecolor{myblue}{RGB}{0,139,139}


%appendix
\usepackage[toc,page]{appendix}

%killing indent
\setlength{\parindent}{0pt}
\usepackage{multicol}
\usepackage{siunitx}
\usepackage{hyperref}


\title{Studying the $Z$ Boson with the ATLAS Detector at the LHC}
\author{Aaron Mielke & Thomas Ackermann}

\begin{document}


\begin{center}
	\makeatletter
	\thispagestyle{empty}
	
	\begin{figure}[H]
	\flushright
	\includegraphics[width=0.35\textwidth]{fig/logo}
	\end{figure}
	
	\vspace{-30mm}
	
	\begin{flushleft}
	\large{\textbf{Fortgeschrittenen Praktikum} \\
		Summer term 2019} \\
	\end{flushleft}
	
	\vspace{5mm}
	
	\rule{\textwidth}{0.2pt}

	\vspace{50mm}
	\Huge\textbf{\@title} \\
	\vspace{10mm}
	\large{\@author} \\
	\normalfont
	
	\vspace{2mm}
	
	\makeatother
\end{center}

\normalsize
\newpage

\tableofcontents

\section{Introduction}

The goal of the lab course was to analyze data from the ATLAS experiment and 
to calculate the mass of the $Z$ Boson.

\subsection{Drell-Yan Process}
A $Z$ Boson can be created during the so called ``Drell-Yan'' Process.
When a quark and an anti-quark collide either a virtual photon or a $Z$ Boson can be produced.

\subsection{ATLAS Detector}
The Detector consists of three main components: inner detector, calorimeters and the muon spectrometer.
The inner detector is mainly used to reconstruct the trajectories of electrically charged particles.



\section{Versuchsdurchführung}




\end{document}

